\documentclass[11pt,a4paper]{article}
\usepackage[margin=26mm,headheight=14pt]{geometry}
\usepackage[T1]{fontenc}
\usepackage{amsmath,amsthm,mathtools}
\usepackage{newtxtext,newtxmath}
\usepackage{microtype,booktabs,array,longtable,enumitem,fancyhdr,xurl}
\usepackage{xcolor}
\usepackage[colorlinks=true,linkcolor=blue!45!black,citecolor=blue!45!black,urlcolor=blue!45!black]{hyperref}
\usepackage[nameinlink,noabbrev]{cleveref}
\numberwithin{equation}{section}
\newtheorem{theorem}{Theorem}[section]
\newtheorem{lemma}[theorem]{Lemma}
\newtheorem{proposition}[theorem]{Proposition}
\newtheorem{corollary}[theorem]{Corollary}
\theoremstyle{definition}\newtheorem{definition}[theorem]{Definition}
\theoremstyle{remark}\newtheorem{remark}[theorem]{Remark}
\newcommand{\Q}{\mathbb Q}\newcommand{\C}{\mathbb C}\newcommand{\R}{\mathbb R}\newcommand{\N}{\mathbb N}
\newcommand{\ee}{\mathrm e}\newcommand{\ii}{\mathrm i}\newcommand{\dd}{\,\mathrm d}
\newcommand{\Li}{\operatorname{Li}}\newcommand{\PV}{\operatorname{PV}}
\newcommand{\Rea}{\operatorname{Re}}\newcommand{\Ima}{\operatorname{Im}}
\newcommand{\B}{\mathcal B}\newcommand{\calF}{\mathcal F}\newcommand{\calL}{\mathcal L}
\newcommand{\coef}[2]{[#1^{#2}]}
\newcommand{\snap}{\texttt{8f5afd4d7565df808359f6a0c22f87f79a119bb9}}
\setlength{\emergencystretch}{2em}
\setlist{itemsep=3pt,topsep=5pt}
\pagestyle{fancy}\fancyhf{}
\fancyhead[L]{\small Order-jet filters for polylogarithmic transseries}
\fancyhead[R]{\small ProveIt research}
\fancyfoot[C]{\thepage}
\hypersetup{pdftitle={Order-Jet Filters for Polylogarithmic Transseries},
 pdfauthor={Research report prepared for Vladimir Reshetnikov},
 pdfsubject={Sharp cancellation costs, Stokes-neutral curves, Stieltjes constants, and digamma inverse transseries}}

\begin{document}
\begin{center}
{\small PROVEIT RESEARCH REPORT \quad\textbar\quad 12 OCTOBER 2026}\par\vspace{8mm}
{\LARGE\bfseries Order-Jet Filters for\par Polylogarithmic Transseries}\par\vspace{5mm}
{\large Sharp cancellation costs, Stokes-neutral curves,\par
and a Stieltjes--digamma inverse}\par\vspace{7mm}
{Prepared for Vladimir Reshetnikov and the ProveIt project}
\end{center}
\vspace{4mm}
\begin{abstract}
We study finite differential operators in the order of
$F_s(x)=-\Li_s(-\ee^x)$, retaining both the Stokes jump and the
exponentially small correction to the arithmetic lateral average.
The classical completion formula turns an order polynomial $\Pi$ into
its values at $-\log n\pm\ii\pi$. This gives a sharp cancellation theorem:
the real polynomials cancelling both data at the first $N$ actions form
exactly the ideal generated by
$\prod_{n=1}^N((u+\log n)^2+\pi^2)$.
We also determine the separate optimal costs, including their change at
integer and half-integer orders, and use Gelfond--Schneider to prove a
nonvanishing rule for every pure integer-order jet.
For each $N\ge1$ a normalized degree-$N+1$ filter admits a unique
convergent moving-order curve on which the entire jump is zero. Its
detuning is explicitly asymptotic to a constant times $\ee^{-Nx}$,
although its physical completion still starts at $\ee^{-x}$.
At order zero, Jonqui\`ere's formula gives all jets in generalized
Stieltjes constants. A positive endpoint-subtracted integral shows that
$H(x)=\partial_sF_s(x)|_{s=0}$ is a real-analytic increasing bijection
$\R\to(0,\infty)$. Its digamma core has an exact ordinary Borel
representation. We prove convergent inverse charts at both endpoints
and compute a model-specific arithmetic-averaging defect:
the inverse-average error is asymptotic to $\pi^2w^2\ee^{-2w}$,
whereas the physical correction to the exact core inverse is only
$(\log2)w\ee^{-2w}$. Full proofs, finite exact regressions, and
independent high-precision integral checks are supplied.
\end{abstract}
\medskip
\noindent\textbf{Status.} Conventional mathematical proofs; not externally
refereed or Lean-formalized. The cancellation classification, moving-order
curve, and assembled order-zero inverse theory are proposed contributions.
Worldwide priority is not established. Classical completion, Gamma/Bell
calculus, and general nonlinear inverse transport are explicitly credited.
\newpage
\tableofcontents
\newpage

\section{The research question and its source boundary}
\label{ojf:sec:context}

\subsection{Three objects that must be kept separate}
A terminating or factorially divergent power--logarithmic expansion does
not determine a unique analytic function. In this report the distinction
can be made exact. For $x>0$ set
\begin{equation}\label{ojf:eq:defs}
 F_s(x)=-\Li_s(-\ee^x),\qquad
 T_s(x)=F_s(-x)=\sum_{n\ge1}\frac{(-1)^{n-1}\ee^{-nx}}{n^s}.
\end{equation}
The subscript is the \emph{polylogarithm order}; the usual normalized
Fermi--Dirac index is $s-1$ \cite{DLMFpoly}.
There are two distinguished lateral realizations $L_s^\pm$ of the
Sommerfeld expansion, satisfying
\begin{equation}\label{ojf:eq:intro-completion}
 L_s^\pm=F_s+\ee^{\mp\ii\pi s}T_s.
\end{equation}
Consequently the jump $L_s^+-L_s^-$ and the physical correction
$(L_s^++L_s^-)/2-F_s$ are different data. Neither should be silently
substituted for the other. We use \emph{arithmetic average}, not
``median,'' for $(L_s^++L_s^-)/2$.

For a real polynomial $\Pi$, applying $\Pi(\partial_s)$ changes the
coefficient of the $n$th exponential into the value of $\Pi$ at
$-\log n\mp\ii\pi$. We ask how many initial actions can be removed by
an order jet of prescribed degree, how this cancellation behaves under
small order detuning, and what analytic inversion does to the remaining
completion. Here an \emph{action $n$} simply means the exponential
$\ee^{-nx}$; this terminology does not add a summability hypothesis.

\subsection{What the repository already proves}
The comparison snapshot is
\begin{center}\small\snap.\end{center}
At that snapshot the package \emph{Resonant Polylogarithm Transseries}
\cite{RepoResonant}, dated 11 October 2026, already proves an explicit
lattice Borel formula, the subtracted Gamma kernel for positive-integer
first order derivatives, and harmonic Bell-polynomial formulas for
all late coefficients. It also proves sharp dilogarithm inverse-radius
results. We use its notation and its completion theorem; we do not
present those results as new. Its Research Question 4 concerns global
higher-jet Borel normal forms, and Question 6 concerns periodic-weight
cancellation. Our order-polynomial filters are a different, precisely
specified cancellation mechanism, not a claimed solution of all the
periodic-weight or higher-kernel questions.

The general nonlinear inverse-transport theorem and the obstruction to
commuting arithmetic averaging with inversion already appear in
\emph{Nonlinear Stokes Transport under Logarithmic Inversion}
\cite{RepoTransport}. The present inverse calculations specialize that
architecture to an exact digamma core. Complete Fermi--Dirac expansions,
including exponentially small terms, are classical; a directly inspected
account is Gil--Segura--Temme, Section 3 and Lemma 3.1 \cite{GST}.
Jonqui\`ere inversion and the Gamma and digamma identities are also
classical \cite{DLMFpoly,DLMFgamma,DLMFpsi}.

The new assertions proposed here are the sharp jet-degree classification,
its moving-order zero curves and nonvanishing consequences, and the
specific endpoint/inverse theory of $H$. Several ingredients are elementary
once the correct exact completion is chosen. Their value lies in a
complete classification with sharp converses and analytic realizations,
not in a claim that polynomial interpolation itself is new.

\subsection{What was not inspected}
The incoming directory listing at the pinned snapshot contained
\path{ProveIt_Lambert_Polylog_Order_Jets_2026-10-12.zip},
\path{ProveIt_Quadratic_Polylog_Boundaries_2026-10-12.zip}, and
\path{Resonant_Lambert_Inverse_Transseries_2026-10-12.zip}.
Their binary contents could not be retrieved through the available
text-only GitHub read route; alternative retrieval attempts did not
supply them. Their names are \emph{not} evidence of their theorem contents.
Thus nonduplication relative to those three archives, and relative to all
unread canonical chapters, is not certified. The source audit records the
actual read ranges and this limitation. No remote repository file was
changed, and no new mathematical error was established in the accepted
statements actually inspected.

\section{Exact lateral completion and analytic continuation}
\label{ojf:sec:completion}

\subsection{The analytic realization used throughout}
We first state the existing completion in the exact normalization needed
below, and give its derivation to fix signs and domains.

\begin{proposition}[Classical completion in the repository normalization]
\label{ojf:prop:completion}
For $x>0$ and $\Rea s>0$, the formal Sommerfeld series is
\begin{equation}\label{ojf:eq:sommerfeld}
 \widehat P_s(x)=\frac{x^s}{\Gamma(s+1)}+
       2\sum_{k\ge1}\frac{\eta(2k)}{\Gamma(s-2k+1)}x^{s-2k},
 \qquad \eta(z)=(1-2^{1-z})\zeta(z).
\end{equation}
In the Borel convention
$\B(\sum_{j\ge1}a_jx^{-j})(\xi)=
 \sum_{j\ge1}a_j\xi^{j-1}/(j-1)!$,
the normalized minor is
\begin{equation}\label{ojf:eq:lattice}
 b_s(\xi)=s\sum_{n\ge1}\frac{(-1)^{n-1}}n
 \left\{(1+\xi/n)^{s-1}-(1-\xi/n)^{s-1}\right\}.
\end{equation}
The upper and lower lateral realizations
\begin{equation}\label{ojf:eq:lateral-borel}
 L_s^\pm(x)=\frac{x^s}{\Gamma(s+1)}
       \left(1+\int_{0}^{\infty\ee^{\pm\ii0}}
                      \ee^{-x\xi}b_s(\xi)\dd\xi\right)
\end{equation}
satisfy \eqref{ojf:eq:intro-completion}. Equivalently, with
$a_\pm=\frac12\pm\frac{\ii x}{2\pi}$,
\begin{equation}\label{ojf:eq:hurwitz-completion}
 L_s^\pm(x)=-\frac{(2\pi)^s\ee^{\mp\ii\pi s/2}}{\Gamma(s)}
                     \zeta(1-s,a_\pm).
\end{equation}
The last formula has a removable singularity at $s=0$ and is entire in
$s$. It supplies the continuation of the lateral realization to every
order. All finite order derivatives may be taken in the exact completion.
\end{proposition}

\begin{proof}
For $\Rea s>0$ the Fermi integral is
\[
 F_s(x)=\frac1{\Gamma(s)}\int_0^\infty
                    \frac{t^{s-1}}{1+\ee^{t-x}}\dd t.
\]
Split at $t=x$ and put $u=|t-x|$. This gives
\begin{equation}\label{ojf:eq:thermal-split}
 F_s(x)=\frac{x^s}{\Gamma(s+1)}+
 \frac1{\Gamma(s)}\int_0^\infty
 \frac{(x+u)^{s-1}-{\bf1}_{u<x}(x-u)^{s-1}}{\ee^u+1}\dd u.
\end{equation}
The binomial expansion of \eqref{ojf:eq:lattice} at zero has coefficient
$2\eta(2k)s^{\underline{2k}}/(2k-1)!$ at $\xi^{2k-1}$.
This proves the formal Borel identity. The lattice sum converges normally
on compact continued sheets away from $\pm1,\pm2,\ldots$: the bracket
in its $n$th term is $O(n^{-1})$ and its prefactor contributes another
$n^{-1}$. The same estimates, with logarithmic factors, hold for finitely
many order derivatives on compact order sets.

On a ray strictly above or below the positive axis the kernel has at most
polynomial--logarithmic growth. One can verify this by splitting the sum
at $n\asymp|\xi|$: the larger $n$ use the preceding difference estimate,
and the finitely many smaller $n$ use the power growth of their binomials.
Thus the tilted Laplace integrals converge. On the boundary their local
powers are integrable when $\Rea s>0$.

In the $n$th integral substitute $\xi=nu/x$. Summing
$\sum_{n\ge1}(-1)^{n-1}\ee^{-nu}=1/(\ee^u+1)$ turns the integral part
of \eqref{ojf:eq:lateral-borel} into
\[
 \frac1{\Gamma(s)}\int_0^\infty
 \frac{(x+u)^{s-1}-(x-u)_{\mp}^{s-1}}{\ee^u+1}\dd u.
\]
The upper path gives argument $-\pi$ to $x-u$ when $u>x$.
Subtracting \eqref{ojf:eq:thermal-split} therefore gives
\[
 -\frac{\ee^{\mp\ii\pi(s-1)}}{\Gamma(s)}
  \int_x^\infty\frac{(u-x)^{s-1}}{\ee^u+1}\dd u
   =\ee^{\mp\ii\pi s}T_s(x).
\]
Endpoint domination by $v^{\epsilon-1}(1+|\log v|^r)$ on compact
sets $\Rea s\ge\epsilon>0$, and exponential domination at infinity,
justify the differentiations and interchanges. Finite Taylor expansion
of the Borel germ with a remainder gives the asymptotic interpretation
of \eqref{ojf:eq:sommerfeld}; no infinite divergent sum was integrated
term by term.

Jonqui\`ere's identity \cite[Eq.~25.12.13]{DLMFpoly}, at $z=-\ee^x$,
gives the lower formula in \eqref{ojf:eq:hurwitz-completion}; conjugating
for real $s$ and then using analytic continuation gives the upper one.
For $a_\pm$ in the right half-plane, $\zeta(1-s,a_\pm)$ has only its
simple pole at $s=0$, cancelled by the simple zero of $1/\Gamma(s)$.
The small-argument series $T_s$ is entire in $s$, normally on compact
sets. Hence both the completion and the physical $F_s$ continue
entirely in the order, and finite order derivatives remain valid.
\end{proof}

\begin{remark}[Continuation is not an unregularized boundary integral]
For $\Rea s\le0$, \eqref{ojf:eq:hurwitz-completion} is the definition of
the continued lateral realization used here. We do not assert that the
boundary powers in \eqref{ojf:eq:lateral-borel} become ordinary locally
integrable functions there. In Section~\ref{ojf:sec:borelzero} we give a
separate, ordinary meromorphic Borel representation for the first jet
at zero. Differentiating the \emph{whole} normalized expression matters:
$\log x$ is not treated as a constant coefficient in an unjustified
ordinary Borel transform.
\end{remark}

\subsection{A polynomial dictionary for the exponential data}
For real $\sigma$ and $\Pi\in\R[u]$, write
\begin{align}
 \calF_{\sigma,\Pi}(x)&=
       \left.\Pi(\partial_s)F_s(x)\right|_{s=\sigma},\label{ojf:eq:observable}\\
 \calL^\pm_{\sigma,\Pi}(x)&=
       \left.\Pi(\partial_s)L_s^\pm(x)\right|_{s=\sigma},\notag\\
 C_{\sigma,\Pi}&=\tfrac12(\calL^+_{\sigma,\Pi}+
                           \calL^-_{\sigma,\Pi})-\calF_{\sigma,\Pi},\notag\\
 J_{\sigma,\Pi}&=\calL^+_{\sigma,\Pi}-\calL^-_{\sigma,\Pi}.\notag
\end{align}
Put $D=\mathrm d/\mathrm du$. The operators
\begin{equation}\label{ojf:eq:operators}
 R_\sigma=\cos(\pi\sigma+\pi D),\qquad
 I_\sigma=\sin(\pi\sigma+\pi D)
\end{equation}
are their Taylor series in the nilpotent operator $D$ on any fixed-degree
polynomial space; all sums terminate.

\begin{theorem}[Exact order-filter dictionary]\label{ojf:thm:dictionary}
Let $u_n=-\log n$ and $q=\ee^{-x}$. Then
\begin{align}
 \calL^\pm_{\sigma,\Pi}-\calF_{\sigma,\Pi}
 &=\ee^{\mp\ii\pi\sigma}
       \sum_{n\ge1}\frac{(-1)^{n-1}}{n^\sigma}
                      \Pi(u_n\mp\ii\pi)q^n,\label{ojf:eq:filter-full}\\
 C_{\sigma,\Pi}
 &=\sum_{n\ge1}\frac{(-1)^{n-1}}{n^\sigma}(R_\sigma\Pi)(u_n)q^n,
                                                        \label{ojf:eq:filter-C}\\
 J_{\sigma,\Pi}
 &=-2\ii\sum_{n\ge1}\frac{(-1)^{n-1}}{n^\sigma}(I_\sigma\Pi)(u_n)q^n.
                                                        \label{ojf:eq:filter-J}
\end{align}
The series converge absolutely for $|q|<1$, normally on compact order
sets. For real $u$,
\begin{equation}\label{ojf:eq:rotation}
 (R_\sigma\Pi)(u)+\ii(I_\sigma\Pi)(u)
      =\ee^{\ii\pi\sigma}\Pi(u+\ii\pi).
\end{equation}
\end{theorem}
\begin{proof}
Every derivative of
$\ee^{\mp\ii\pi s}n^{-s}$ multiplies it by
$-\log n\mp\ii\pi$. Apply the finite operator to
\eqref{ojf:eq:intro-completion}. A polynomial in $\log n$ times
$n^{-\sigma}$ has subexponential growth in $n$, which proves the stated
normal convergence. Taylor's formula for a polynomial gives
$\Pi(u+\ii\pi)=\exp(\ii\pi D)\Pi(u)$ and hence
\eqref{ojf:eq:rotation}. Taking the arithmetic average and difference
proves the two remaining formulas, with the sign in the jump fixed by
upper minus lower.
\end{proof}

\section{Sharp cancellation costs and a complete joint ideal}
\label{ojf:sec:filters}

\subsection{What a degree can and cannot remove}
For a nonzero power series in $q$ with no constant term, let $\nu$ be its
first nonzero exponent; let $\nu(0)=\infty$.

\begin{theorem}[Sharp separate degree costs]\label{ojf:thm:cost}
Fix $\Pi\ne0$ of degree $d$ and real $\sigma$.
If $\sin\pi\sigma\ne0$, then $\nu(J_{\sigma,\Pi})\le d+1$;
if $\sin\pi\sigma=0$ and $d\ge1$, then
$\nu(J_{\sigma,\Pi})\le d$.
Similarly, if $\cos\pi\sigma\ne0$, then
$\nu(C_{\sigma,\Pi})\le d+1$;
if $\cos\pi\sigma=0$ and $d\ge1$, then
$\nu(C_{\sigma,\Pi})\le d$.
Every displayed bound is attained by a polynomial of the indicated degree.
The exceptional constant polynomials have zero jump at integer order,
and zero arithmetic completion at half-integer order.
\end{theorem}

\begin{proof}
The constant term of $I_\sigma$ as a polynomial in $D$ is
$\sin\pi\sigma$. When this is nonzero, $I_\sigma$ is an invertible
triangular operator on degree-at-most-$d$ polynomials and preserves the
degree. When it is zero,
\[
 I_\sigma=(\pi\cos\pi\sigma)D+O(D^3),
\]
so it has kernel exactly the constants and reduces every positive degree
by exactly one. The analogous assertions for $R_\sigma$ follow from
its constant and linear terms, $\cos\pi\sigma$ and
$-\pi\sin\pi\sigma$.

A nonzero polynomial of degree $k$ cannot vanish at all of the $k+1$
distinct real numbers $u_1,\ldots,u_{k+1}$. The coefficients multiplying
these evaluations in \eqref{ojf:eq:filter-C}--\eqref{ojf:eq:filter-J}
never vanish. This proves the upper bounds.

To attain a bound for a degree-preserving operator, take its inverse
image of $\prod_{n=1}^d(u+\log n)$. For a degree-lowering operator,
its map from polynomials of degree at most $d$, modulo constants, to
polynomials of degree at most $d-1$ is an isomorphism. Take the inverse
image of $\prod_{n=1}^{d-1}(u+\log n)$. The next node is not a root,
so the bound is attained exactly. The assertions for constants follow
directly from \eqref{ojf:eq:operators}.
\end{proof}

Thus, for $d\ge1$, the optimal first surviving actions are
\begin{center}
\begin{tabular}{@{}lll@{}}\toprule
Order center & Jump & Arithmetic completion\\\midrule
$\sigma\in\mathbb Z$ & $d$ & $d+1$\\
$\sigma\in\mathbb Z+\tfrac12$ & $d+1$ & $d$\\
Neither of the above & $d+1$ & $d+1$\\\bottomrule
\end{tabular}
\end{center}
The two columns are separately optimized; they need not be realized by
the same polynomial. In particular, a degree statement about one datum
cannot be reused for their simultaneous cancellation.

For example, at integer order the operators are, up to $(-1)^\sigma$,
$\cos(\pi D)$ and $\sin(\pi D)$. A polynomial that delays the completion
to action $d+1$ is
\begin{equation}\label{ojf:eq:cos-extremizer}
 \Pi^{C}_d(u)=\sec(\pi D)\prod_{n=1}^d(u+\log n).
\end{equation}
The power series for $\sec$ terminates on the polynomial. A primitive
followed by the finite operator $\pi D/\sin(\pi D)$ similarly constructs
a jump extremizer. These formulas are exact algebra, not asymptotic
infinite operator expansions.

\subsection{Simultaneous cancellation is a principal ideal}

\begin{theorem}[Joint cancellation ideal and optimal real jet degree]
\label{ojf:thm:joint}
Let $E$ be any finite set of positive integers. For every real order
center $\sigma$, the real polynomials $\Pi$ for which both the jump and
the arithmetic completion vanish at every action in $E$ form exactly
\begin{equation}\label{ojf:eq:joint-ideal}
 \left(\prod_{n\in E}\bigl((u+\log n)^2+\pi^2\bigr)\right)\R[u].
\end{equation}
Consequently, for $E=\{1,\ldots,N\}$ the minimum nonzero degree is $2N$,
and the unique monic minimizer is
\begin{equation}\label{ojf:eq:joint-extremizer}
 Q_N^{\mathrm{joint}}(u)=
       \prod_{n=1}^N\bigl((u+\log n)^2+\pi^2\bigr).
\end{equation}
At action $N+1$, at least one of its jump and completion coefficients
is nonzero. More generally a degree-$d$ nonzero jet cannot make both
data zero at all of the first $\lfloor d/2\rfloor+1$ actions.
\end{theorem}

\begin{proof}
By \eqref{ojf:eq:rotation}, both real components at the node $u_n$
vanish if and only if $\Pi(u_n+\ii\pi)=0$. Since $\Pi$ is real,
this is equivalent to zeros at both $u_n+\ii\pi$ and $u_n-\ii\pi$,
or to divisibility by $(u-u_n)^2+\pi^2$. These real quadratic factors
are pairwise coprime: all their complex roots are distinct. Divisibility
by each is therefore equivalent to divisibility by their product.
This proves the ideal identity and the degree assertions.

For the monic minimizer at the next node, putting
$\ell_j=\log(j/(N+1))$ gives
\begin{equation}\label{ojf:eq:next-joint}
 Q_N^{\mathrm{joint}}(-\log(N+1)+\ii\pi)
      =\prod_{j=1}^N\ell_j(\ell_j+2\ii\pi)\ne0.
\end{equation}
After rotation by $\ee^{\ii\pi\sigma}$ this still has at least one
nonzero real component. This proves the final assertion. We do not
claim that both components must separately be nonzero in this formula.
\end{proof}

\begin{corollary}[First-order detuning robustness]\label{ojf:cor:robustness}
For a fixed real polynomial $\Pi$ and a finite action set $E$, the
following requirements are equivalent: the jump coefficients at every
$n\in E$ vanish together with their first order-center derivatives at
one real center $\sigma_0$; both the jump and completion coefficients
vanish there; and the full lateral completion coefficients vanish at
those actions for \emph{every} order center. Their common solution ideal
is \eqref{ojf:eq:joint-ideal}. In particular robustness through first
order in detuning costs exactly $2N$ for the first $N$ actions.
\end{corollary}
\begin{proof}
Up to a nonzero fixed scalar, the $n$th jump coefficient is
$n^{-\sigma}(I_\sigma\Pi)(u_n)$. Its derivative is
\[
 n^{-\sigma}\{-\log n\,(I_\sigma\Pi)(u_n)
                      +\pi(R_\sigma\Pi)(u_n)\}.
\]
When the jump coefficient is zero, its derivative is zero exactly when
the completion coefficient is zero. The joint theorem then forces
$\Pi(u_n\pm\ii\pi)=0$, independently of the order center. The converse
is immediate from the exact dictionary.
\end{proof}

\begin{corollary}[Finite reconstruction from exponential data]
\label{ojf:cor:reconstruction}
Fix $N$ and real $\sigma$. A polynomial of degree less than $2N$ is
uniquely determined by its jump and arithmetic-completion coefficients
at the first $N$ actions. Every assignment of those $2N$ real normalized
data is attained by exactly one such polynomial.
\end{corollary}
\begin{proof}
Divide out the nonzero scalar coefficients in
\eqref{ojf:eq:filter-C}--\eqref{ojf:eq:filter-J} and undo the rotation
in \eqref{ojf:eq:rotation}. This specifies $\Pi$ at the $N$ upper nodes
and, by conjugacy, at the $N$ lower nodes. Ordinary interpolation at
$2N$ distinct complex points gives a unique degree-less-than-$2N$
polynomial. Its conjugate-reflection polynomial has the same values,
so uniqueness makes its coefficients real. Equivalently, injectivity
follows from \cref{ojf:thm:joint} and equal finite dimensions give
surjectivity.
\end{proof}

\paragraph{The first joint filter.}
For $N=1$, the operator is $\partial_s^2+\pi^2$. It removes the
entire first action, not just its imaginary part. At integer $m$ the
next completion coefficient is
$(-1)^{m+1}(\log2)^2/2^m$ and the next jump coefficient is
$-4\pi\ii(-1)^m\log2/2^m$. This is a simple exact test of the
real-versus-imaginary distinction.

\subsection{An explicit finite-tail estimate}
For numerical evaluation the exact coefficient dictionary also supplies
an elementary bound. Let $\Pi(u)=\sum_{j=0}^dp_ju^j$ and
$K_\Pi=\sum_{j=0}^d|p_j|(1+\pi)^j$. For real $\sigma\ge0$,
$0<q<1$, and integer $M\ge0$, each full lateral completion tail obeys
\begin{equation}\label{ojf:eq:tailbound}
 \left|\sum_{n>M}\frac{(-1)^{n-1}}{n^\sigma}
                  \Pi(-\log n\pm\ii\pi)q^n\right|
 \le K_\Pi(M+1)^{d-\sigma}q^{M+1}
                  \frac{E_d(q)}{(1-q)^{d+1}},
\end{equation}
where $E_d$ is the polynomial defined by
$\sum_{k\ge0}(k+1)^dq^k=E_d(q)/(1-q)^{d+1}$, with $E_0=1$.
Indeed $|\Pi(-\log n\pm\ii\pi)|\le K_\Pi n^d$ for $n\ge1$,
$n=M+1+k\le(M+1)(k+1)$, and $n^{-\sigma}\le(M+1)^{-\sigma}$.
This proves the bound. The jump tail has at most twice the right side;
the arithmetic-completion tail has at most the right side. This is a
finite evaluation safeguard, not the basis of the cancellation theorem.

\section{Pure-jet rigidity and exact Stokes-neutral order curves}
\label{ojf:sec:neutral}

\subsection{A transcendence input rules out accidental zeros}

\begin{lemma}\label{ojf:lem:log-transcendence}
For every integer $n\ge2$, $\log n/\pi$ is transcendental.
\end{lemma}
\begin{proof}
Use the classical Gelfond--Schneider theorem in its complex form: if
$\alpha\ne0,1$ and $\beta\notin\Q$ are algebraic, every fixed value of
$\alpha^\beta$ is transcendental \cite[Theorem 3.1]{GS}.
If $a=\log n/\pi$ were algebraic, $\beta=-\ii a$ would be nonreal
algebraic, hence irrational. With $\log(-1)=\ii\pi$,
$(-1)^{-\ii a}=\exp(\pi a)=n$, a contradiction.
\end{proof}

\begin{theorem}[Parity and nonvanishing for every pure jet]
\label{ojf:thm:parity}
Fix an integer order $m$ and $r\ge1$, and take $\Pi(u)=u^r$.
At every action $n\ge2$, both the jump coefficient and the
arithmetic-completion coefficient are nonzero. Their first surviving
actions are exactly
\begin{equation}\label{ojf:eq:parity-table}
 (\nu(J_{m,u^r}),\nu(C_{m,u^r}))=
 \begin{cases}(1,2),&r\ \text{odd},\\(2,1),&r\ \text{even}.
 \end{cases}
\end{equation}
For $r=0$, the jump is identically zero and the completion starts at
one. At half-integer centers the two columns of
\eqref{ojf:eq:parity-table} are exchanged.
\end{theorem}
\begin{proof}
At $n=1$, the relevant complex value is $(\ii\pi)^r$; its real and
imaginary parts alternate with the parity of $r$. At $n\ge2$, divide
$(-\log n+\ii\pi)^r$ by $\pi^r$. Its real and imaginary parts are
nonzero polynomials in $a=\log n/\pi$ with integer coefficients:
the real polynomial has leading term $(-a)^r$, and the imaginary one
has leading term $r(-a)^{r-1}$. Neither can vanish at the transcendental
number in \cref{ojf:lem:log-transcendence}. At integer order the rotation
is just $(-1)^m$; at half-integer order it is $\pm\ii$, which exchanges
the two components. The constant case follows from the completion formula.
\end{proof}

\begin{corollary}[Rigidity for algebraic normalized coefficients]
\label{ojf:cor:algebraic-rigidity}
The nonvanishing at every action $n\ge2$ extends from pure monomials to
all nonconstant filters
$\Pi(u)=c\,p(u/\pi)$, where $c\in\R\setminus\{0\}$ and $p$ has
real algebraic coefficients. At integer or half-integer centers no such
filter can cancel even one of the two data at action two. In particular
arbitrarily high degree in this coefficient class cannot reproduce the
deep logarithmic-node filters of \cref{ojf:thm:cost}.
\end{corollary}
\begin{proof}
The real and imaginary parts of $p(-a+\ii)$ are nonzero polynomials in
$a$ with real algebraic coefficients: their degrees are respectively
$\deg p$ and $\deg p-1$. A number transcendental over $\Q$ cannot be
a root of a nonzero polynomial over the algebraic numbers. Apply
\cref{ojf:lem:log-transcendence} and then the integer or half-integer
rotation. The common nonzero real scale $c$ does not affect vanishing.
\end{proof}

This use of Gelfond--Schneider proves an all-order nonvanishing statement;
a finite table of numerical coefficients would not prove it. Conversely,
the freely chosen real coefficients in \cref{ojf:thm:joint} deliberately
permit cancellation. There is no contradiction: a pure monomial in the
order derivative is a much smaller class of filters.

\subsection{A normalized sharp jump filter}
For $N\ge1$ define
\begin{equation}\label{ojf:eq:QN}
 Q_N(u)=\prod_{j=1}^N(u+\log j).
\end{equation}
There is a unique real polynomial $\Pi_N$ with
\begin{equation}\label{ojf:eq:PiN}
 \sin(\pi D)\Pi_N=Q_N,\qquad
      (\cos(\pi D)\Pi_N)(0)=1.
\end{equation}
It has degree $N+1$. Existence and uniqueness follow from the
isomorphism modulo constants used in \cref{ojf:thm:cost}; the second
condition fixes that constant. One constructive formula is
\[
 \Pi_N=\frac1\pi\frac{\pi D}{\sin(\pi D)}
               \int_0^u Q_N(v)\dd v+c_N,
\]
where $c_N$ enforces the normalization. All operators terminate.
Writing $R_N=\cos(\pi D)\Pi_N$, the first two examples are
\begin{align}
 \Pi_1(u)&=\frac{u^2}{2\pi}+1+\frac\pi2,
 &R_1(u)&=1+\frac{u^2}{2\pi},\label{ojf:eq:Pi1}\\
 \Pi_2(u)&=\frac{u^3}{3\pi}+
       \frac{\log2}{2\pi}u^2+\frac\pi3u+1+\frac{\pi\log2}{2},
 &R_2(u)&=1+\frac{u^3}{3\pi}+
       \frac{\log2}{2\pi}u^2-\frac{2\pi}{3}u.
                                                        \label{ojf:eq:Pi2}
\end{align}
The normalization is intentionally about the \emph{completion}, not
about $\Pi_N(0)$.

\begin{theorem}[Convergent Stokes-neutral curve and its crossover]
\label{ojf:thm:neutral}
Fix an integer $m\ge0$ and $N\ge1$. For sufficiently small $q$, there
is a unique holomorphic function $\epsilon_N(q)$ with
$\epsilon_N(0)=0$ for which
\begin{equation}\label{ojf:eq:neutral-identity}
 J_{m+\epsilon_N(q),\Pi_N}(-\log q)=0
\end{equation}
when $q>0$. For complex $q$, this identity refers to the convergent
$q$-series of the jump. The function has real Taylor coefficients and
\begin{equation}\label{ojf:eq:neutral-leading}
 \epsilon_N(q)=-\frac{\kappa_{m,N}}\pi q^N+O(q^{N+1}),\qquad
 \kappa_{m,N}=\frac1{(N+1)^m}
                 \prod_{j=1}^N\log\frac{N+1}{j}>0.
\end{equation}
The associated completion does \emph{not} disappear:
\begin{equation}\label{ojf:eq:neutral-completion}
 C_{m+\epsilon_N(q),\Pi_N}(-\log q)
                  =(-1)^m q+O(q^2).
\end{equation}
More generally, with $\epsilon=\lambda q^N$,
\begin{equation}\label{ojf:eq:crossover}
 \frac{J_{m+\lambda q^N,\Pi_N}(-\log q)}
                   {-2\ii(-1)^m q^{N+1}}
       =\pi\lambda+\kappa_{m,N}+O(q),
\end{equation}
uniformly for $\lambda$ in any fixed compact subset of $\C$.
At $m=0$ these are analytically continued lateral realizations; no
unregularized order-zero boundary integral is assumed.
\end{theorem}

\begin{proof}
Remove the factor $-2\ii(-1)^m q$ from the jump and define
\begin{align}
 \Psi(\epsilon,q)=\sum_{n\ge1}(-1)^{n-1}n^{-m-\epsilon}q^{n-1}
 \{\sin(\pi\epsilon)R_N(u_n)+
                  \cos(\pi\epsilon)Q_N(u_n)\}.\label{ojf:eq:Psi}
\end{align}
The series is holomorphic for $|q|<1$ and $\epsilon$ in a neighborhood
of zero: on compact sets its coefficients have at most a fixed power of
$n$ times a fixed power of $\log n$. Its $n=1$ term is precisely
$\sin(\pi\epsilon)$, since $Q_N(0)=0$ and $R_N(0)=1$.
Consequently
\[
 \Psi(0,0)=0,\qquad \partial_\epsilon\Psi(0,0)=\pi.
\]
The holomorphic implicit-function theorem gives the unique curve.
Conjugation and uniqueness imply real Taylor coefficients.

At $\epsilon=0$, the first $N$ coefficients vanish. At $n=N+1$,
\[
 (-1)^N(N+1)^{-m}Q_N(-\log(N+1))=\kappa_{m,N}.
\]
Thus $\Psi(0,q)=\kappa_{m,N}q^N+O(q^{N+1})$. Substitution in the
implicit equation, whose $\epsilon$ derivative is $\pi+O(q)$, gives
\eqref{ojf:eq:neutral-leading}. Taylor expansion with
$\epsilon=\lambda q^N$ proves \eqref{ojf:eq:crossover}; higher sine
terms and mixed terms are of sufficiently high order, and holomorphy
gives uniformity on compact $\lambda$ sets.

The corresponding completion has normalized coefficient
\[
 \cos(\pi\epsilon)R_N(u_n)-\sin(\pi\epsilon)Q_N(u_n).
\]
Its first term is $(-1)^m q\cos(\pi\epsilon)$. The rest is $O(q^2)$
uniformly for small $\epsilon$, and the implicit curve tends to zero.
This proves \eqref{ojf:eq:neutral-completion}.
\end{proof}

\begin{remark}[What the curve cancels]
For a \emph{fixed} order, the sharp filters remove individual initial
action coefficients. On the moving-order curve, coefficients at different
actions compensate after $\epsilon$ is substituted as a function of $q$.
The result is an exact identity of analytic functions, not a statement
that every fixed-order Borel singularity has disappeared. It also gives
a precise example in which zero total lateral jump does not identify the
physical function with the common lateral realization.
\end{remark}

\section{All order-zero jets in Stieltjes and integral coordinates}
\label{ojf:sec:stieltjes}

\subsection{A complete finite formula in generalized Stieltjes constants}
For $\Rea a>0$ use the Hurwitz--Stieltjes convention \cite{DLMFhurwitz}
\begin{equation}\label{ojf:eq:stieltjes}
 \zeta(1+z,a)=\frac1z+
       \sum_{n\ge0}\frac{(-1)^n\gamma_n(a)}{n!}z^n.
\end{equation}
In particular $\gamma_0(a)=-\psi(a)$.
Define the complete exponential Bell polynomials by
\begin{equation}\label{ojf:eq:bell}
 \exp\left(\sum_{j\ge1}b_j\frac{z^j}{j!}\right)
       =\sum_{r\ge0}\B_r(b_1,\ldots,b_r)\frac{z^r}{r!},
 \qquad\B_0=1.
\end{equation}
The following theorem specializes classical Jonqui\`ere inversion to a
resonance excluded from the positive-integer Gamma-kernel statement of
\cite{RepoResonant}.

\begin{theorem}[Stieltjes coordinates for every zero-order lateral jet]
\label{ojf:thm:stieltjes}
Let $x>0$, $a_\pm=\frac12\pm\frac{\ii x}{2\pi}$, and
\begin{align}
 d_1^\pm&=\gamma+\log(2\pi)\mp\frac{\ii\pi}{2},\label{ojf:eq:dj}\\
 d_j^\pm&=(-1)^{j+1}(j-1)!\zeta(j),\qquad j\ge2.\notag
\end{align}
For $r\ge1$ put
\begin{equation}\label{ojf:eq:stieltjes-jet}
 Z_r^\pm(x)=\B_r(d^\pm)
       -r\sum_{n=0}^{r-1}\binom{r-1}{n}
                    \gamma_n(a_\pm)\B_{r-1-n}(d^\pm),
 \qquad Z_0^\pm=1.
\end{equation}
Then $Z_r^\pm=\partial_s^rL_s^\pm|_{s=0}$ and
\begin{equation}\label{ojf:eq:physical-zero-jets}
 \left.\partial_s^rF_s(x)\right|_{s=0}
  =Z_r^\pm(x)-\sum_{n\ge1}(-1)^{n-1}
                     (-\log n\mp\ii\pi)^r\ee^{-nx}.
\end{equation}
The formulas hold for every finite $r$ and locally uniformly for $x>0$.
\end{theorem}

\begin{proof}
The reciprocal-Gamma expansion \cite[\S5.7]{DLMFgamma} gives
\[
 \frac1{\Gamma(s)}=s\exp\left(\gamma s+
           \sum_{j\ge2}\frac{(-1)^{j+1}\zeta(j)}j s^j\right),
 \qquad |s|<1.
\]
In \eqref{ojf:eq:stieltjes} substitute $z=-s$. Combining it with
\eqref{ojf:eq:hurwitz-completion} yields the exact local generating
identity
\begin{equation}\label{ojf:eq:stieltjes-generating}
 L_s^\pm(x)=
 \exp\left(\sum_{j\ge1}d_j^\pm\frac{s^j}{j!}\right)
       \left(1-\sum_{n\ge0}\gamma_n(a_\pm)\frac{s^{n+1}}{n!}\right).
\end{equation}
In particular the removable value at zero is $1$, not zero. Multiplying
Taylor series and taking $r!$ times the coefficient gives
\eqref{ojf:eq:stieltjes-jet}. Applying the exact completion formula at
zero gives \eqref{ojf:eq:physical-zero-jets}. The local Taylor identities
and the normally convergent exponential series justify every finite
coefficient extraction and the claimed local uniformity.
\end{proof}

For example,
\[
 Z_1^\pm=d_1^\pm-\gamma_0(a_\pm),\qquad
 Z_2^\pm=(d_1^\pm)^2-\zeta(2)
                   -2d_1^\pm\gamma_0(a_\pm)-2\gamma_1(a_\pm).
\]
The order-zero second physical derivative is therefore an explicit
combination of $\gamma_1$ at a conjugate pair of points, digamma values,
and a convergent log-weighted exponential series. There is no finite-part
ambiguity in these identities: the regularization is fixed by the
holomorphic product in \eqref{ojf:eq:hurwitz-completion}.

\subsection{Independent endpoint-subtracted integral coordinates}
A second exact representation is useful both for proof and for independent
numerical evaluation. It avoids differentiating a numerically delicate
polylogarithm continuation at an integer order.

\begin{theorem}[Positive regularization and all zero-order moments]
\label{ojf:thm:regularization}
For any real $x$ let $q=\ee^{-x}$ and $f_0=(1+q)^{-1}$. Define
\begin{equation}\label{ojf:eq:moments}
 M_j(x)=\frac1{1+q}\int_0^\infty
       \frac{(1-\ee^{-t})(\log t)^j}{t(1+\ee^{t-x})}\dd t,
       \qquad j\ge0.
\end{equation}
All these integrals are absolutely convergent. With
$g_1=\gamma$ and $g_j=(-1)^{j+1}(j-1)!\zeta(j)$ for $j\ge2$,
\begin{equation}\label{ojf:eq:moment-jets}
 F_0(x)=f_0,\qquad
 \left.\partial_s^rF_s(x)\right|_{s=0}
   =r\sum_{j=0}^{r-1}\binom{r-1}{j}\B_{r-1-j}(g)M_j(x)
       \quad(r\ge1).
\end{equation}
In particular $H(x):=\partial_sF_s(x)|_{s=0}=M_0(x)>0$.
\end{theorem}

\begin{proof}
Let $f_x(t)=(1+\ee^{t-x})^{-1}$. The subtraction has the exact
positive form
\begin{equation}\label{ojf:eq:positive-subtraction}
 f_x(t)-f_0\ee^{-t}
  =\frac{1-\ee^{-t}}{(1+\ee^{-x})(1+\ee^{t-x})}.
\end{equation}
For $\Rea s>0$, subtract the elementary Gamma integral to get
\begin{equation}\label{ojf:eq:regularized-F}
 F_s(x)=f_0+\frac1{\Gamma(s)}\int_0^\infty
                          t^{s-1}(f_x(t)-f_0\ee^{-t})\dd t.
\end{equation}
At zero the subtracted numerator is $O(t)$, so the remaining integral
is holomorphic for $\Rea s>-1$. At infinity it has exponential decay
on compact real $x$ sets. Multiplication by $1/\Gamma(s)$ continues
\eqref{ojf:eq:regularized-F} to a neighborhood of zero.
Its $j$th integral derivative at zero is $M_j(x)$; the majorants
$t^{-1+\epsilon}(1+|\log t|^j)$ after subtraction near zero, with
$\epsilon>0$, and exponential majorants at infinity justify this.
The reciprocal-Gamma Bell expansion used above now gives
\eqref{ojf:eq:moment-jets}. Positivity of the integrand for $j=0$ proves
the final assertion.
\end{proof}

\section{The order-zero digamma core and its exact Borel kernel}
\label{ojf:sec:borelzero}

\subsection{An exact physical identity}
Put
\begin{align}
 A(x)&=\gamma+\log(2\pi)+
              \Rea\psi\left(\tfrac12+\frac{\ii x}{2\pi}\right),
                   &&x>0,\label{ojf:eq:A-real}\\
 B(z)&=\sum_{n\ge2}(-1)^n(\log n)z^n
          =(\log2)z^2-(\log3)z^3+(\log4)z^4-\cdots,
                   &&|z|<1.\label{ojf:eq:B}
\end{align}
The letter $B$ here denotes a convergent analytic function, not the
Borel transform $\B$.

\begin{proposition}[Digamma core and both lateral first jets]
\label{ojf:prop:digamma}
For $x>0$,
\begin{align}
 H(x)&=A(x)-B(\ee^{-x}),\label{ojf:eq:H-A}\\
 Z_1^\pm(x)&=A(x)\mp\ii\pi\frac{\ee^{-x}}{1+\ee^{-x}}.
                                                        \label{ojf:eq:Z1}
\end{align}
Thus the first zero-order jet has a jump beginning at action one,
while its physical correction to the arithmetic average begins at
action two.
\end{proposition}
\begin{proof}
By \cref{ojf:thm:stieltjes},
$Z_1^+=\gamma+\log(2\pi)-\ii\pi/2+\psi(a_+)$.
The digamma reflection formula \cite[\S5.5]{DLMFgamma} gives
\[
 \psi(1-z)-\psi(z)=\pi\cot(\pi z),\qquad
 \Ima\psi\left(\tfrac12+\frac{\ii x}{2\pi}\right)
                =\frac\pi2\tanh(x/2).
\]
Hence the imaginary part of $Z_1^+$ is $-\pi/(\ee^x+1)$, and
its real part is $A$. The lower formula follows by conjugation.
In \eqref{ojf:eq:physical-zero-jets} with $r=1$, the imaginary tail
cancels exactly, while the real tail is $B(\ee^{-x})$.
\end{proof}

\begin{corollary}[Symmetry and the global positive core branch]
\label{ojf:cor:core-global}
The even function $A$ satisfies
\[
 H(x)+H(-x)=A(x),\qquad A(0)=\log(\pi/2),
 \qquad H(0)=\tfrac12\log(\pi/2).
\]
It is strictly increasing for $x>0$ and maps that half-line onto
$(\log(\pi/2),\infty)$.
\end{corollary}
\begin{proof}
For $x>0$, $H(-x)=B(\ee^{-x})$, giving the symmetry from
\eqref{ojf:eq:H-A}; continuity and evenness give the remaining real
points. The classical value $\psi(1/2)=-\gamma-2\log2$ gives $A(0)$.
For $v=x/(2\pi)>0$, the convergent trigamma series gives
\[
 A'(x)=-\frac1{2\pi}\Ima\sum_{k\ge0}\frac1{(k+1/2+\ii v)^2}>0,
\]
because each summand has strictly negative imaginary part. The range
then follows from the logarithmic asymptotic proved below. This also
shows that the exact positive core inverse used later is unique wherever
its target exceeds $\log(\pi/2)$.
\end{proof}

The single meromorphic expression, holomorphic near large real points, is
\begin{equation}\label{ojf:eq:A-holomorphic}
 A(z)=\gamma+\log(2\pi)+\frac12
    \left\{\psi\left(\tfrac12+\frac{\ii z}{2\pi}\right)
          +\psi\left(\tfrac12-\frac{\ii z}{2\pi}\right)\right\}.
\end{equation}
It is meromorphic in $z$, with possible poles on the imaginary axis,
and agrees with \eqref{ojf:eq:A-real} on the positive real axis. In
particular it is holomorphic on every disk of a fixed radius around a
sufficiently large positive real $w$. We use this expression, not the
nonholomorphic operation $\Rea$, in complex inverse equations.

\subsection{A meromorphic minor at the previously excluded order}

\begin{theorem}[Ordinary Borel reconstruction at order zero]
\label{ojf:thm:borelzero}
The formal algebraic tail of $A(x)-(\log x+\gamma)$ is
\begin{equation}\label{ojf:eq:A-asymptotic}
 A(x)\sim\log x+\gamma-
          2\sum_{k\ge1}\eta(2k)(2k-1)!x^{-2k}.
\end{equation}
Its ordinary Borel minor is the meromorphic function
\begin{equation}\label{ojf:eq:b0}
 b_0^{\rm jet}(\xi)=\frac1\xi-\frac\pi{\sin\pi\xi},
\end{equation}
with a removable zero at the origin and residue $(-1)^{n-1}$ at
$\xi=n\ge1$. For $x>0$,
\begin{align}
 \log x+\gamma+\int_\pm\ee^{-x\xi}b_0^{\rm jet}(\xi)\dd\xi
     &=Z_1^\pm(x),\label{ojf:eq:b0-lateral}\\
 A(x)&=\log x+\gamma+
          \PV\int_0^\infty\ee^{-x\xi}b_0^{\rm jet}(\xi)\dd\xi.
                                                        \label{ojf:eq:b0-PV}
\end{align}
The principal value is the arithmetic average of the lateral integrals,
or the sum of symmetric local principal values at the positive integers.
\end{theorem}

\begin{proof}
The classical generator
$\pi\xi/\sin\pi\xi=1+2\sum_{k\ge1}\eta(2k)\xi^{2k}$
shows that
\[
 b_0^{\rm jet}(\xi)=-2\sum_{k\ge1}\eta(2k)\xi^{2k-1}.
\]
The ordinary Borel convention gives exactly the coefficients in
\eqref{ojf:eq:A-asymptotic}. The residue statement follows from
$(\sin\pi\xi)'|_{\xi=n}=\pi(-1)^n$.

For an independent proof of the exact reconstruction, initially take
$\Rea x>0$ and $\Ima x<0$. Rotate an upper tilted Laplace path to
$\xi=\ii t$. There are no poles inside the rotation sector. On its
large arcs, the cosecant has exponential decay away from the real
axis and the $1/\xi$ term is suppressed by $\ee^{-x\xi}$; the negative
imaginary part of $x$ ensures decay even at the final imaginary ray.
The removable origin gives no small-circle contribution. Thus
\begin{align*}
 \int_+\ee^{-x\xi}b_0^{\rm jet}(\xi)\dd\xi
 &=\int_0^\infty \ee^{-\ii xt}
                   \left(\frac1t-\frac\pi{\sinh\pi t}\right)\dd t\\
 &=\int_0^\infty
       \left\{\frac{\ee^{-\ii xu/(2\pi)}}u
          -\frac{\ee^{-(1/2+\ii x/(2\pi))u}}{1-\ee^{-u}}\right\}\dd u.
\end{align*}
The standard digamma integral, combined with the Frullani difference,
is \cite[\S5.9]{DLMFpsi}
\[
 \psi(a)=\int_0^\infty
          \left(\frac{\ee^{-u}}u-\frac{\ee^{-au}}{1-\ee^{-u}}\right)\dd u,
 \quad \Rea a>0,
\]
where cancellation at zero is taken inside the integrand. Consequently
the preceding integral is
\begin{equation}\label{ojf:eq:b0-evaluation}
 \psi\left(\tfrac12+\frac{\ii x}{2\pi}\right)
                 -\log x+\log(2\pi)-\frac{\ii\pi}{2}.
\end{equation}
Both sides continue to a neighborhood of the positive real axis. A fixed
upper ray of angle less than $\pi/2$ has an absolutely convergent
Laplace integral there, so \eqref{ojf:eq:b0-evaluation} proves the upper
formula for real positive $x$. Conjugation proves the lower formula.
Its value is precisely \eqref{ojf:eq:Z1}.

On the real boundary the only singularities are simple poles. Upper
minus lower indentation contributes minus $2\pi\ii$ times their
residues, and the weighted residue sum is absolutely convergent:
\[
 -2\pi\ii\sum_{n\ge1}(-1)^{n-1}\ee^{-nx}
          =-2\pi\ii\frac{\ee^{-x}}{1+\ee^{-x}}.
\]
At infinity one may take boundaries midway between the integer poles;
exponential decay controls the remaining contour and principal-value
tails. Averaging proves \eqref{ojf:eq:b0-PV}.
Finally, a finite Taylor expansion of the germ on a tilted ray, with a
Laplace remainder, gives \eqref{ojf:eq:A-asymptotic} to each finite order.
\end{proof}

The kernel \eqref{ojf:eq:b0} is the natural $m=0$ extension of the
subtracted Gamma kernel in \cite[label
\texttt{thm:Gamma-kernel}]{RepoResonant}. The extension has been proved
here with a separate contour calculation, rather than by asserting
that all hypotheses of the $m\ge1$ theorem persist unchanged.

\begin{corollary}[Derivative scales for the exact core]
\label{ojf:cor:core-derivatives}
As $w\to+\infty$,
\begin{align}
 A(w)&=\log w+\gamma-\frac{\pi^2}{6w^2}+O(w^{-4}),\label{ojf:eq:A-two}\\
 A'(w)&=\frac1w+O(w^{-3}),\qquad
 A''(w)=-\frac1{w^2}+O(w^{-4}).\label{ojf:eq:A-derivatives}
\end{align}
The estimates are holomorphically uniform in a fixed closed subsector
about the positive real axis, away from the origin. In particular,
$A'(w)>0$ for all sufficiently large real $w$.
\end{corollary}
\begin{proof}
Use the first term of \eqref{ojf:eq:A-asymptotic}. Uniformity follows
from the tilted Laplace proof, or the standard sectorial digamma
expansion \cite[Eq.~5.11.2]{DLMFasymptotic}, applied to the conjugate-pair
formula \eqref{ojf:eq:A-holomorphic}. Cauchy estimates on disks of radius
a fixed positive multiple of $w$, contained in such a subsector, justify
differentiation of each finite asymptotic remainder and give
\eqref{ojf:eq:A-derivatives}. This uses holomorphic remainder control,
not mere differentiation of a real-axis big-$O$ statement.
\end{proof}

\section{A global inverse and convergent charts at both endpoints}
\label{ojf:sec:inverse}

\subsection{The real inverse exists globally}

\begin{theorem}[Global order-zero inverse]\label{ojf:thm:global}
The function $H(x)=\partial_sF_s(x)|_{s=0}$ is a real-analytic strictly
increasing bijection from $\R$ to $(0,\infty)$. Its inverse $X(y)$ is
real analytic on $(0,\infty)$. Its two endpoint scales are
\begin{equation}\label{ojf:eq:H-endpoints}
 H(x)\sim(\log2)\ee^{2x}\quad(x\to-\infty),\qquad
 H(x)\sim\log x\quad(x\to+\infty).
\end{equation}
\end{theorem}
\begin{proof}
Use the positive integral $H=M_0$ from
\cref{ojf:thm:regularization}. Both factors
$(1+\ee^{-x})^{-1}$ and $(1+\ee^{t-x})^{-1}$ increase strictly with
$x$ for every $t>0$. Differentiation on compact $x$ intervals is
justified by the same endpoint and exponential majorants, and the
integrand derivative is strictly positive. Hence $H'(x)>0$.
Local holomorphic continuation of the integral in $x$ in a strip of
sufficiently small width gives real analyticity.

For $x<0$, the original small-argument polylogarithm series and normal
order differentiation give $H(x)=B(\ee^x)$, so
$H(x)\sim(\log2)\ee^{2x}\to0$. For large positive $x$,
\eqref{ojf:eq:H-A} and \eqref{ojf:eq:A-two} give
$H(x)=\log x+\gamma+O(x^{-2})\to\infty$.
Continuity and strict increase prove the range assertion, and the
analytic inverse-function theorem applies everywhere because $H'>0$.
\end{proof}

\subsection{The small-target logarithmic--square-root chart}
Write $a=\log2$ and
\begin{equation}\label{ojf:eq:small-b}
 b(z)=\frac{B(z)}{az^2}
       =1-\frac{\log3}{a}z+\frac{\log4}{a}z^2-\cdots,
 \qquad b(0)=1.
\end{equation}
The square-root parameter for a small target is
$t=(y/a)^{1/2}>0$.

\begin{theorem}[All coefficients and an explicit convergent endpoint disk]
\label{ojf:thm:small-inverse}
For $0<t<\sqrt3/40$, the global inverse has the convergent expansion
\begin{equation}\label{ojf:eq:small-inverse}
 X(at^2)=\log t+\sum_{N\ge1}c_Nt^N,
 \qquad c_N=\frac1N\coef{z}{N}b(z)^{-N/2}.
\end{equation}
The analytic correction $X(at^2)-\log t$ extends holomorphically to
$|t|<\sqrt3/40$. In particular,
\begin{equation}\label{ojf:eq:small-first}
 X(y)=\frac12\log\frac ya+
     \frac{\log3}{2a}\left(\frac ya\right)^{1/2}
     +\left(\frac{(\log3)^2}{2a^2}-1\right)\frac ya
     +O(y^{3/2}).
\end{equation}
The radius given here is a sufficient certificate, not an exact radius.
\end{theorem}

\begin{proof}
Let $r=1/20$. Since $\log n\le n$ for $n\ge1$,
\[
 \sup_{|z|\le r}|b(z)-1|
 \le\frac1a\sum_{n\ge3}n r^{n-2}
 =\frac{r(3-2r)}{a(1-r)^2}
 =\frac{58}{361a}<\frac14.
\]
For the strict last inequality, $\log2>2/3$ follows, for example, from
$\log2=2\sum_{k\ge0}(1/3)^{2k+1}/(2k+1)$; then
$58/(361a)<87/361<1/4$.
There is a holomorphic square root of $b$ on the disk, equal to one
at zero, and $|\sqrt{b(z)}|>\sqrt3/2$ there. The function
$f(z)=z\sqrt{b(z)}$ has exactly the simple zero $0$, and on $|z|=r$
its modulus exceeds $\sqrt3/40$.

For $|t|<\sqrt3/40$, Rouch\'e's theorem gives exactly one root,
counting multiplicity, of $f(z)-t$ inside $|z|<r$. That root is simple,
so it defines a holomorphic function $z(t)$ with $z(0)=0$ and
$z'(0)=1$. Uniqueness glues the local roots throughout the disk.
The quotient $z(t)/t$ extends holomorphically and has no zeros;
its logarithm with value zero at $t=0$ therefore exists on the disk.
For positive small $t$, the root is real and positive and
$H(\log z(t))=B(z(t))=at^2$. The global uniqueness theorem identifies
$\log z(t)$ with $X(at^2)$.

For the coefficient formula, set $\Phi(z)=b(z)^{-1/2}$. The equation
is $z=t\Phi(z)$ and $\log(z/t)=\log\Phi(z)$.
Lagrange--B\"urmann inversion gives, for $N\ge1$,
\begin{align*}
 \coef{t}{N}\log(z(t)/t)
 &=\frac1N\coef{z}{N-1}(\log\Phi)'\Phi^N\\
 &=\frac1{N^2}\coef{z}{N-1}(\Phi^N)'
  =\frac1N\coef{z}{N}\Phi^N.
\end{align*}
This proves \eqref{ojf:eq:small-inverse}. Expanding the first two
coefficients and using $\log4=2\log2$ proves
\eqref{ojf:eq:small-first}.
\end{proof}

\subsection{The large-target exact-core chart}
For large $w>0$, take the target $y=A(w)$, so $w$ is the exact
inverse of the algebraic digamma core rather than an approximation to
that inverse. Put
\begin{equation}\label{ojf:eq:large-notations}
 q=\ee^{-w},\qquad a_w=A'(w),\qquad b_w=A''(w),\qquad
 c_+=-\ii\pi,\quad c_-=\ii\pi.
\end{equation}
Let $G_\pm(y)$ denote the inverse branches near $w$ of the lateral
functions $Z_1^\pm(x)$ from \eqref{ojf:eq:Z1}.

\begin{theorem}[Convergent inverse sectors and the exact second sector]
\label{ojf:thm:large-inverse}
There are constants $c>0$ and $w_0>0$ such that for each $w>w_0$ the
three inverse displacements are holomorphic for $|q|<c/w$ when
$q$ is treated as an independent variable. At the physical substitution
$q=\ee^{-w}$, their convergent expansions give
\begin{align}
 X(A(w))&=w+\frac{\log2}{a_w}q^2+O(w^3q^3),
                                            \label{ojf:eq:physical-inverse}\\
 G_\pm(A(w))&=w-\frac{c_\pm}{a_w}q+
 \left(-\frac{b_wc_\pm^2}{2a_w^3}
       -\frac{c_\pm^2}{a_w^2}+\frac{c_\pm}{a_w}\right)q^2
                   +O(w^3q^3).              \label{ojf:eq:lateral-inverse}
\end{align}
The remainder constants can be chosen independently of large $w$.
The physical branch in \eqref{ojf:eq:physical-inverse} is the global
inverse of \cref{ojf:thm:global}.
\end{theorem}

\begin{proof}
For independent $q$, the defining equations for the displacements are
\begin{align}
 A(w+\delta)-A(w)-B(q\ee^{-\delta})&=0,
                                             \label{ojf:eq:inv-physical-eq}\\
 A(w+\delta)-A(w)+c_\pm
       \frac{q\ee^{-\delta}}{1+q\ee^{-\delta}}&=0.
                                             \label{ojf:eq:inv-lateral-eq}
\end{align}
Fix a small radius $\rho>0$, for example $\rho=1/4$. The uniform
estimates in \cref{ojf:cor:core-derivatives} imply
\[
 A(w+\delta)-A(w)=\frac\delta w(1+O(w^{-1}))
                    \qquad (|\delta|\le\rho).
\]
For sufficiently large $w$, the core difference has exactly one zero
in the disk and has boundary modulus at least $\rho/(2w)$; this follows
by comparison with $\delta/w$.

On the same disk, the lateral perturbation is bounded by
$\pi |q|\ee^\rho/(1-|q|\ee^\rho)$.
Choose $c>0$ small enough that this is less than $\rho/(2w)$ whenever
$|q|\le c/w$, for all sufficiently large $w$.
The physical perturbation $B(q\ee^{-\delta})$ is $O(|q|^2)$ uniformly
on the disk and satisfies the same strict inequality, after increasing
$w_0$ if necessary. Rouch\'e gives exactly one simple zero in each case.
The implicit-function theorem and uniqueness give holomorphic roots on
the parameter disk, bounded by $\rho$.
Cauchy's estimate bounds their degree-$n$ coefficients by
$\rho(w/c)^n$. Hence on $|q|\le c/(2w)$ the remainder after degree two
is at most
\[
 \rho\frac{(w|q|/c)^3}{1-w|q|/c}=O(w^3|q|^3).
\]
The physical value $q=\ee^{-w}$ lies in this region for all large $w$.
This proves actual sector convergence and a uniform remainder; it is
not only a formal reversion computation.

For the coefficients, write $\delta=d_1q+d_2q^2+O(q^3)$ and use
\[
 A(w+\delta)-A(w)=a_w\delta+\tfrac12b_w\delta^2+O(\delta^3),
 \qquad
 \frac{q\ee^{-\delta}}{1+q\ee^{-\delta}}
       =q-(d_1+1)q^2+O(q^3).
\]
The first coefficient in the lateral equation gives $d_1=-c/a_w$.
The second gives
$a_wd_2+\tfrac12b_wd_1^2-cd_1-c=0$, which is precisely
\eqref{ojf:eq:lateral-inverse}. In the physical equation,
$B(q\ee^{-\delta})=(\log2)q^2+O(q^3)$, so $d_1=0$ and
$d_2=\log2/a_w$.
For the physical real $q$, conjugation and uniqueness make the physical
root real. Its defining equation is $H(w+\delta)=A(w)$, and the global
monotonicity of $H$ identifies it with $X(A(w))$.
\end{proof}

\subsection{A larger inverse error than the physical completion}

\begin{corollary}[Digamma-core arithmetic-averaging defect]
\label{ojf:cor:defect}
With the preceding exact-core normalization,
\begin{align}
 \frac{G_+(A(w))+G_-(A(w))}{2}-X(A(w))
 &=\left(\frac{\pi^2}{a_w^2}
       +\frac{\pi^2b_w}{2a_w^3}-\frac{\log2}{a_w}\right)\ee^{-2w}
             +O(w^3\ee^{-3w})                         \label{ojf:eq:defect-exact}\\
 &\sim\pi^2w^2\ee^{-2w}.                              \label{ojf:eq:defect-leading}
\end{align}
By contrast,
\begin{equation}\label{ojf:eq:physical-leading}
 X(A(w))-w\sim(\log2)w\ee^{-2w},
 \qquad
 G_+(A(w))-G_-(A(w))\sim2\pi\ii w\ee^{-w}.
\end{equation}
In particular the inverse-average error divided by the physical
core-inverse correction tends to infinity, with equivalent
$\pi^2w/\log2$.
\end{corollary}
\begin{proof}
Average \eqref{ojf:eq:lateral-inverse}, using
$c_++c_-=0$ and $c_\pm^2=-\pi^2$, then subtract
\eqref{ojf:eq:physical-inverse}. This gives
\eqref{ojf:eq:defect-exact}. Since
$a_w=w^{-1}(1+O(w^{-2}))$ and
$b_w=-w^{-2}(1+O(w^{-2}))$, the coefficient of $\ee^{-2w}$ in that
formula is $\pi^2w^2+O(w)$.
The uniform remainder is smaller by a factor $O(w\ee^{-w})$.
The same estimates give the first equivalent in
\eqref{ojf:eq:physical-leading}. Subtracting the two lateral
expansions gives its second equivalent. Division proves the ratio.
\end{proof}

This is a model-specific quantitative consequence of the general
nonlinear averaging mechanism already proved in \cite{RepoTransport}
and restated in \cite{RepoResonant}. What is added here is the explicit
core, the exact second-sector coefficient, a uniform convergent chart,
and the separation by an unbounded factor from the physical correction.
It is \emph{not} a counterexample to a correctly formulated nonlinear
median summation. Nor does it allege that either cited repository article
confused arithmetic averaging with such a summation.

\paragraph{In the target variable.}
Let $v=\ee^{y-\gamma}$. Inverting \eqref{ojf:eq:A-two} gives
$w=v+\pi^2/(6v)+O(v^{-3})$. Therefore the physical inverse correction
and the inverse-average error have the respective equivalents
\begin{equation}\label{ojf:eq:double-exponential}
 (\log2)\ee^{y-\gamma}\exp(-2\ee^{y-\gamma}),\qquad
 \pi^2\ee^{2(y-\gamma)}\exp(-2\ee^{y-\gamma}).
\end{equation}
These are doubly exponential scales in $y$. Using the exact $w$ for
finite computations remains preferable: an algebraic approximation to
the core must not be confused with its much smaller exponential
correction.

\section{Verification, source audit, and integration safeguards}
\label{ojf:sec:verification}

\subsection{Retained finite evidence}
The accompanying program \path{code/verify.py} runs without network
access. The retained run used Python 3.13.5, SymPy 1.14.0, and mpmath
1.3.0 at 70 decimal digits. It passed \textbf{165 exact assertions} and
\textbf{46 numerical assertions}, with no numerical group skipped.
The exact checks include finite shift identities, operator inversions
and ranks through degree ten, joint-root and Vandermonde checks through
six pairs of nodes, normalized logarithmic-node filters through $N=3$,
and the second-order inverse coefficient calculation. Rational-node
checks test universal polynomial algebra; actual logarithmic nodes are
used separately where stated. None of these finite checks is a proof of
the all-degree theorems.

For the Stieltjes formulas, the program extracts the coefficients of the
\emph{regularized} function $\zeta(1-s,a)+1/s$ by Cauchy sampling, at
both 64 and 96 points, and compares the reconstructed first four
physical jets at $x=1.3,3$ with the independent integral in
\cref{ojf:thm:regularization}. This reference avoids the small-step
polylogarithm-order differentiation issue recorded in
\cite{RepoResonant}. Two tilted-ray Borel integrals are compared with
the digamma formula. Additional checks cover three noninteger-order
completion identities, six moving-order roots, five pairs of large
inverse residuals, and three small-inverse residuals plus their first
two coefficients. The largest absolute discrepancy among the 46
retained numerical assertions was below $2\cdot10^{-68}$.
These are high-precision floating-point agreements, not interval bounds.

The asymptotic tables are diagnostics, not additional certified
assertions about limits. Table~\ref{ojf:tab:neutral} shows the normalized
order detuning and the still-present completion on the neutral curve.
Here $m=1$, so the signed leading completion is $-q$.
\begin{table}[htbp]\centering\small
\caption{Moving-order diagnostics, rounded from the retained 70-digit run.}
\label{ojf:tab:neutral}
\begin{tabular}{@{}rrrr@{}}\toprule
$N$ & $q$ & $\epsilon_N/(-\kappa_{1,N}q^N/\pi)$ & $C/(-q)$\\\midrule
1 & $10^{-2}$ & 0.9956011681 & 0.9946589499\\
1 & $10^{-3}$ & 0.9995582649 & 0.9994621827\\
2 & $10^{-2}$ & 0.9962146385 & 0.9877618322\\
2 & $10^{-3}$ & 0.9996172905 & 0.9987664010\\
3 & $10^{-2}$ & 0.9888879416 & 0.9895194578\\
3 & $10^{-3}$ & 0.9988740742 & 0.9989479273\\\bottomrule
\end{tabular}
\end{table}

\begin{table}[htbp]\centering\small
\caption{Large-core inverse diagnostics. The last column compares the
actual inverse-average error with the full second-sector term in
\eqref{ojf:eq:defect-exact}, not just its leading equivalent.}
\label{ojf:tab:inverse}
\begin{tabular}{@{}rrrr@{}}\toprule
$w$ & $(X-w)/(w\ee^{-2w})$ & Defect$/ (w^2\ee^{-2w})$
 & Defect/second sector\\\midrule
5  & 0.5853514479 & 5.9934294047 & 0.9633235152\\
8  & 0.6485114953 & 7.9349324606 & 0.9988202952\\
12 & 0.6756899051 & 8.9089971386 & 0.9999811029\\
20 & 0.6872810572 & 9.4220413126 & 0.9999999938\\
35 & 0.6912692028 & 9.6554457036 & 1.0000000000\\\midrule
Limit & $\log2$ & $\pi^2$ & $1$\\\bottomrule
\end{tabular}
\end{table}
The last displayed value is rounded; it does not mean exact equality.
For the small-endpoint chart the error after six terms is about
$3.90\cdot10^{-17}$ at $t=10^{-2}$ and $3.88\cdot10^{-24}$ at
$t=10^{-3}$, consistent with a seventh-order remainder. Full values and
residuals, rather than only these rounded summaries, are retained in
\path{data/}.

\subsection{Corrections and qualifications that belong with this package}
No newly verified mathematical error is attributed to the inspected
repository sources. The following are precise integration safeguards,
not invented corrections to unseen articles.

\paragraph{Do not discard a zero-times-pole product at order zero.}
The factor $1/\Gamma(s)$ vanishes at $s=0$, but the Hurwitz zeta factor
in \eqref{ojf:eq:hurwitz-completion} has a pole there. Their product
has the nonzero value $L_0^\pm=1$. Evaluating the factors separately
would lose all the Stieltjes-jet information.

\paragraph{Do not identify zero jump with physical reconstruction.}
The exact neutral curve in \cref{ojf:thm:neutral} has zero jump and a
nonzero completion starting at action one. Conversely, cancelling both
fixed-order data costs twice as many prescribed roots, as
\cref{ojf:thm:joint} proves. Neither assertion by itself removes all
singularities on all continued Borel sheets.

\paragraph{Keep differentiation and continuation hypotheses.}
The boundary integral for the normalized lattice kernel was proved in
$\Rea s>0$. At zero we used analytic continuation and then a separately
proved meromorphic kernel. A bare substitution into a nonintegrable
boundary power is not an alternative proof. Order derivatives also act
on $x^s/\Gamma(s+1)$, not just on an isolated Borel minor.

\paragraph{Keep the exact inverse core and the meaning of averaging.}
The physical correction to $w=A^{-1}(y)$ is exponentially small.
Replacing $w$ by a low-order algebraic approximation generally creates
a larger error. The arithmetic average of lateral inverses differs from
the physical inverse by the explicit term
\eqref{ojf:eq:defect-exact}; no claim is made here about constructing
an algebra-compatible median summation.

\paragraph{Do not promote a sufficient disk to an exact radius.}
The small-target radius $\sqrt3/40$ and the large-core parameter radius
$c/w$ are proved sufficient. They are not claimed optimal. Likewise,
real monotonicity of $H$ does not prove complex univalence of its global
continuation.

\subsection{Proposed integration and claim boundary}
The package should be placed in a new directory such as
\begin{quote}\small\ttfamily
Analysis/Transseries/docs/series-and-transseries/\\
Polylogarithm\_Order\_Jet\_Filters/
\end{quote}
with a backlink from the polylogarithm reports. The companion
\path{integration/INTEGRATION.md} identifies proposed insertion points
and result labels. The existing resonant-polylogarithm package should
remain the credited source of the positive-integer Gamma/Bell formulas,
and the existing transport package should remain the credited source
of the general inverse mechanism.

The proofs in this article are complete within their stated hypotheses.
The source comparison is bounded, and the proposed novelty of the
assembled results has not been independently refereed or cleared against
every incoming draft. The package contains neither a Lean development nor
an interval-certified numerical implementation. Its reproducible finite
checks supplement, rather than replace, the mathematical proofs.

\section{Further research questions}
\label{ojf:sec:future}

\paragraph{1. Global Borel normal forms for higher order jets.}
The exact dictionary controls the entire exponential jump, but not an
explicit normal form for every local singular germ of a higher-jet Borel
transform. Develop those germs in a form compatible with the Stieltjes
coordinates and determine which fixed-order filters remove a singularity,
which only suppress a jump coefficient, and what occurs on later sheets.
This continues Research Question 4 of \cite{RepoResonant}; it is not
resolved by the finite root-counting theorem alone.

\paragraph{2. Conditioning and optimal normalization of deep filters.}
The minimal degrees are exact, but the coefficients of the extremizing
polynomials can be large. Determine norms and condition numbers for
interpolation on $-\log n\pm\ii\pi$, and optimize cancellation under a
fixed coefficient norm or a restriction to a prescribed coefficient
field. Exact degree optimality is not numerical optimality.

\paragraph{3. A uniform theory with growing cancellation depth.}
The neutral-curve theorem fixes $N$ before $q\to0$. Determine the
largest ranges $N=N(x)\to\infty$ for which the crossover and the implicit
curve have uniform errors, and the competing growth of $\Pi_N$,
$\kappa_{m,N}$, and the convergence radius. The individual fixed-$N$
proofs do not justify this joint limit.

\paragraph{4. Neutral curves under nonlinear inverse transport.}
When the order is selected as a function of $q=\ee^{-x}$, inversion
must include the induced order dependence in its derivative. Construct
inverse charts along the exact neutral curve, and determine whether
jump cancellation remains useful after the physical completion and
this additional chain rule are both retained. A forward identity alone
does not solve that nonlinear problem.

\paragraph{5. The exact small-target inverse radius of $H$.}
Locate the closest complex critical value or singular value of
$z\sqrt{B(z)/(\log2\,z^2)}$. Is the dominant obstruction a simple fold,
the logarithmic boundary at $z=-1$, or another continued singularity?
This would replace the explicit sufficient disk in
\cref{ojf:thm:small-inverse} by a sharp theorem, with coefficient
asymptotics rather than only local convergence.

\paragraph{6. Higher zero-order jets and their inverse scales.}
The leading algebraic terms of higher jets are Bell polynomials in
$\log x+\gamma$ and zeta values. Classify their eventual monotonicity,
critical points, and inverse branches, and determine which parity-induced
completion actions become doubly exponential sectors after inversion.
The global positivity argument here is special to the first jet and
cannot simply be asserted for the higher logarithmic moments.

\paragraph{7. Combining periodic weights and order filters.}
The periodic-weight dictionary in \cite{RepoResonant} can be combined
with order differentiation. Determine optimal designs when both the
period and jet degree are constrained, keeping the zero Fourier mode.
In particular, quantify the tradeoff between removing residue classes
and removing individual logarithmic spectral nodes; the two forms of
arithmetic freedom should not be conflated.

\paragraph{8. Formal proofs and certified numerical continuation.}
The principal-ideal theorem and finite operator inversions are practical
first formalization targets over real polynomials. The transcendence
corollary can import Gelfond--Schneider rather than recreate it. For the
analytic part, certify the regularized integrals, tilted Borel contours,
and inverse disks with interval enclosures and explicit tail bounds.
The existing numerical scripts deliberately do not claim these certificates.

\appendix
\section{Reproduction and theorem map}
\label{ojf:app:reproduce}

The source is self-contained and uses standard TeX packages; no external
images or custom font files are required. From the package root:
\begin{verbatim}
python -m pip install -r requirements.txt
python code/verify.py --output-dir replay
make pdf
\end{verbatim}
A single verification group can be rerun with, for example,
\begin{verbatim}
python code/verify.py --group stieltjes --output-dir replay-stieltjes
\end{verbatim}
The program refuses to overwrite existing output unless
\texttt{--overwrite} is explicitly provided. Numerical precision is
selectable with \texttt{--dps}, with a minimum of 50 digits. The retained
run uses 70. Each group records its versions, precision, finite assertion
counts, and numerical differences. The supplied build-validation record
reports the actual PDF checks; it should not be confused with a formal
mathematical certificate.

\small
\begin{longtable}{@{}>{\raggedright\arraybackslash}p{.29\textwidth}
 >{\raggedright\arraybackslash}p{.24\textwidth}
 >{\raggedright\arraybackslash}p{.40\textwidth}@{}}
\toprule Result & Status & Main dependency\\\midrule\endhead
\Cref{ojf:prop:completion} & Credited background & Classical Fermi split,
Jonqui\`ere identity; explicit repository realization.\\[4pt]
\Cref{ojf:thm:dictionary} & Derived exact dictionary & Finite order
operator and normally convergent tail.\\[4pt]
\Cref{ojf:thm:cost,ojf:thm:joint} & Proposed sharp classification &
Finite differential operators, polynomial roots, conjugate interpolation.\\[4pt]
\Cref{ojf:thm:parity} & Derived all-order rigidity &
Gelfond--Schneider, imported as a classical theorem.\\[4pt]
\Cref{ojf:thm:neutral} & Proposed analytic cancellation theorem &
Sharp normalized filter and holomorphic implicit functions.\\[4pt]
\Cref{ojf:thm:stieltjes,ojf:thm:regularization} & Exact order-zero synthesis &
Hurwitz Laurent expansion, reciprocal Gamma, endpoint subtraction.\\[4pt]
\Cref{ojf:thm:borelzero} & Proved zero-order kernel extension &
Digamma integral and a specified contour rotation.\\[4pt]
\Cref{ojf:thm:global,ojf:thm:small-inverse} & Proposed endpoint package &
Positive integral, convergent series, Lagrange inversion, Rouch\'e.\\[4pt]
\Cref{ojf:thm:large-inverse,ojf:cor:defect} & Model-specific transport &
Exact digamma core; known general inverse architecture credited.\\\bottomrule
\end{longtable}
\normalsize

\begin{thebibliography}{20}
\small\raggedright
\bibitem{DLMFpoly}
NIST Digital Library of Mathematical Functions,
\emph{\S25.12, Polylogarithms}, especially Eqs.~25.12.11,
25.12.13, and 25.12.16.
\url{https://dlmf.nist.gov/25.12}. Accessed 12 October 2026.

\bibitem{DLMFgamma}
NIST Digital Library of Mathematical Functions,
\emph{\S5.5, Functional Relations; \S5.7, Series Expansions}.
\url{https://dlmf.nist.gov/5.5}; \url{https://dlmf.nist.gov/5.7}.
Accessed 12 October 2026.

\bibitem{DLMFpsi}
NIST Digital Library of Mathematical Functions,
\emph{\S5.9, Integral Representations}, digamma integrals.
\url{https://dlmf.nist.gov/5.9}. Accessed 12 October 2026.

\bibitem{DLMFasymptotic}
NIST Digital Library of Mathematical Functions,
\emph{\S5.11, Asymptotic Expansions}, especially Eq.~5.11.2.
\url{https://dlmf.nist.gov/5.11}. Accessed 12 October 2026.

\bibitem{DLMFhurwitz}
NIST Digital Library of Mathematical Functions,
\emph{\S25.11, Hurwitz Zeta Function}.
\url{https://dlmf.nist.gov/25.11}. Accessed 12 October 2026.

\bibitem{GST}
A. Gil, J. Segura, and N. M. Temme,
Complete asymptotic expansions for the relativistic Fermi--Dirac integral,
\emph{Applied Mathematics and Computation} \textbf{412} (2022), 126618.
\url{https://doi.org/10.1016/j.amc.2021.126618}.
Inspected preprint: arXiv:2108.11210v2, Section 3 and Lemma 3.1,
\url{https://arxiv.org/html/2108.11210v2}.

\bibitem{GS}
M. Karatarakis and F. Wiedijk,
\emph{A formalization of the Gelfond--Schneider theorem},
arXiv:2603.24823v1 (25 March 2026), especially Theorem 3.1 and its
complex-logarithm convention.
\url{https://arxiv.org/html/2603.24823v1}.
The classical theorem is imported here; no formalization of the present
article is claimed.

\bibitem{RepoResonant}
ProveIt research package,
\emph{Resonant Polylogarithm Transseries: Gamma Borel Kernels,
Cyclotomic Order Jets, and Sharp Dilogarithm Reversion},
11 October 2026.
\path{Analysis/Transseries/docs/series-and-transseries/Resonant_Polylogarithm_Transseries/article.tex}.
Snapshot \snap; source blob
\texttt{d54638582ab147eb4a48a846c3bdc371643aaf39}.
Theorems identified in the source by labels
\texttt{thm:lattice}, \texttt{thm:Gamma-kernel}, \texttt{thm:Bell},
and \texttt{thm:periodic}; inversion and audit sections were also inspected
as recorded in the accompanying source audit.

\bibitem{RepoTransport}
ProveIt research package,
\emph{Nonlinear Stokes Transport under Logarithmic Inversion},
\path{Analysis/Transseries/docs/series-and-transseries/Nonlinear_Stokes_Transport_Logarithmic_Inversion/article.tex}.
Snapshot \snap; source blob
\texttt{1e8b24325554d1f1c94c29f21b794ba62feff755}.
The general contour inverse coefficients and the explicit distinction
between arithmetic averaging and nonlinear-compatible median summation
are prior repository material.

\bibitem{RepoInventories}
ProveIt Contributors,
\emph{Polylogarithms project README; Series and transseries inventory;
Incoming-directory listing}, inspected at snapshot \snap.
\path{Analysis/Polylogarithms/README.md};
\path{Analysis/Transseries/docs/series-and-transseries/README.md};
\path{docs/incoming/}.
Repository: \url{https://github.com/VladimirReshetnikov/ProveIt}.
The listing alone does not establish the contents of its binary archives.
\end{thebibliography}
\end{document}
